% Auto-generated from raw.txt
% Usage:
%   \vocabentry{word}{/ipa/ (pos)}{definition}{example}{note}

\vocabentry{Nutrient}
{/ˈnuːtriənt/ (n)}
{A substance that provides nourishment essential for growth and the maintenance of life.}
{Plants draw essential nutrients from the soil to grow.}
{Đây là một danh từ dùng để chỉ chất dinh dưỡng nói chung.}

\vocabentry{Carbohydrate}
{/ˌkɑːrboʊˈhaɪdreɪt/ (n)}
{A large group of organic compounds occurring in foods and living tissues and including sugars, starch, and cellulose.}
{Bread and potatoes are high in carbohydrates, providing quick energy for the body.}
{Đây là một danh từ dùng để chỉ nhóm chất bột đường.}

\vocabentry{Metabolism}
{/məˈtæbəlɪzəm/ (n)}
{The chemical processes that occur within a living organism in order to maintain life.}
{Regular exercise can help to increase your body's metabolism.}
{Đây là một danh từ dùng để chỉ quá trình trao đổi chất.}

\vocabentry{Preservative}
{/prɪˈzɜːrvətɪv/ (n)}
{A substance used to preserve foodstuffs, wood, or other materials against decay.}
{Many processed foods contain artificial preservatives to extend their shelf life.}
{Đây là một danh từ dùng để chỉ chất bảo quản thực phẩm.}

\vocabentry{Calorie}
{/ˈkæləri/ (n)}
{The amount of energy needed to raise the temperature of 1 gram of water by 1 degree Celsius.}
{To lose weight, you should consume fewer calories than you burn each day.}
{Đây là một danh từ dùng để chỉ đơn vị năng lượng calo.}

\vocabentry{Protein}
{/ˈproʊtiːn/ (n)}
{Large group of organic compounds which are essential for the body's structure and function.}
{Lean meat, beans, and eggs are excellent sources of protein for muscle repair.}
{Đây là một danh từ dùng để chỉ chất đạm.}

\vocabentry{Obesity}
{/oʊˈbiːsəti/ (n)}
{The state of being grossly fat or overweight.}
{Obesity is a major health concern linked to heart disease and diabetes.}
{Đây là một danh từ dùng để chỉ bệnh béo phì.}

\vocabentry{Malnutrition}
{/ˌmælnuˈtrɪʃn/ (n)}
{Lack of proper nutrition, caused by not having enough to eat, not eating enough of the right things, or being unable to use the food that one does eat.}
{Eradicating malnutrition is a primary goal of many global health organizations.}
{Danh từ chỉ sự suy dinh dưỡng.}

\vocabentry{Additive}
{/ˈædətɪv/ (n)}
{A substance added to something in small quantities, typically to improve or preserve it.}
{Some food additives have been linked to behavioral problems in children.}
{Đây là một danh từ dùng để chỉ phụ gia thực phẩm.}

\vocabentry{Cholesterol}
{/kəˈlestərɔːl/ (n)}
{A compound of the sterol type found in most body tissues.}
{High levels of cholesterol in the blood can increase the risk of a heart attack.}
{Đây là một danh từ dùng để chỉ chất béo gây xơ vữa động mạch.}

\vocabentry{Vegetarian}
{/ˌvedʒəˈteriən/ (n/adj)}
{A person who does not eat meat; relating to this practice.}
{She decided to become a vegetarian for ethical and health reasons.}
{Đây là một danh từ/tính từ dùng để chỉ người ăn chay hoặc thực phẩm chay.}

\vocabentry{Vegan}
{/ˈviːɡən/ (n/adj)}
{A person who does not eat or use animal products.}
{Vegan diets avoid all animal-derived ingredients, including honey and eggs.}
{Đây là một danh từ/tính từ dùng để chỉ người ăn chay trường/thuần chay.}

\vocabentry{Organic}
{/ɔːrˈɡænɪk/ (adj)}
{Produced or involving production without the use of chemical fertilizers or pesticides.}
{Organic vegetables are often more expensive but contain fewer chemical residues.}
{Đây là một tính từ dùng để chỉ thực phẩm hữu cơ.}

\vocabentry{Digest}
{/daɪˈdʒest/ (v)}
{Break down food in the alimentary canal into substances that can be absorbed by the body.}
{It takes several hours for the body to fully digest a heavy meal.}
{Đây là một động từ dùng để chỉ việc tiêu hóa thức ăn.}

\vocabentry{Indigestion}
{/ˌɪndɪˈdʒestʃən/ (n)}
{Pain or discomfort in the stomach associated with difficulty in digesting food.}
{Eating too quickly can often lead to a painful bout of indigestion.}
{Đây là một danh từ dùng để chỉ chứng khó tiêu.}

\vocabentry{Fiber}
{/ˈfaɪbər/ (n)}
{Dietary material containing substances such as cellulose that are resistant to the action of digestive enzymes.}
{A diet high in fiber is essential for maintaining a healthy digestive system.}
{Đây là một danh từ dùng để chỉ chất xơ.}

\vocabentry{Mineral}
{/ˈmɪnərəl/ (n)}
{An inorganic substance that is essential for the healthy functioning of the body.}
{Calcium is an important mineral for building strong bones and teeth.}
{Đây là một danh từ dùng để chỉ chất khoáng.}

\vocabentry{Vitamin}
{/ˈvaɪtəmɪn/ (n)}
{Any of a group of organic compounds which are essential for normal growth and nutrition.}
{Oranges are a great source of Vitamin C, which helps boost the immune system.}
{Đây là một danh từ dùng để chỉ vitamin.}

\vocabentry{Gluten}
{/ˈɡluːtn/ (n)}
{A substance present in cereal grains, especially wheat, that is responsible for the elastic texture of dough.}
{Many people now choose to follow a gluten-free diet due to intolerance.}
{Đây là một danh từ dùng để chỉ một loại protein trong ngũ cốc.}

\vocabentry{Allergy}
{/ˈælərdʒi/ (n)}
{A damaging immune response by the body to a substance.}
{Peanut allergies can be life-threatening and require immediate medical attention.}
{Đây là một danh từ dùng để chỉ sự dị ứng thực phẩm.}

\vocabentry{Cuisine}
{/kwɪˈziːn/ (n)}
{A style or method of cooking, especially as characteristic of a particular country or region.}
{Italian cuisine is famous worldwide for its pasta and pizza dishes.}
{Đây là một danh từ dùng để chỉ ẩm thực.}

\vocabentry{Delicacy}
{/ˈdelɪkəsi/ (n)}
{A choice or expensive food.}
{Caviar is considered a luxury delicacy in many parts of the world.}
{Đây là một danh từ dùng để chỉ món cao lương mỹ vị hoặc đặc sản quý hiếm.}

\vocabentry{Fermentation}
{/ˌfɜːrmenˈteɪʃn/ (n)}
{The chemical breakdown of a substance by bacteria, yeasts, or other microorganisms.}
{Yogurt is produced through the fermentation of milk by specific bacteria.}
{Đây là một danh từ dùng để chỉ sự lên men.}

\vocabentry{Garnish}
{/ˈɡɑːrnɪʃ/ (n/v)}
{A small amount of food used to decorate other food; to decorate food.}
{The chef used a sprig of parsley as a garnish for the grilled fish.}
{Đây là một danh từ/động từ dùng để chỉ đồ trang trí hoặc việc trang trí món ăn.}

\vocabentry{Ingredient}
{/ɪnˈɡriːdiənt/ (n)}
{Any of the foods or substances that are combined to make a particular dish.}
{Fresh herbs are a key ingredient in many traditional Mediterranean recipes.}
{Đây là một danh từ dùng để chỉ nguyên liệu nấu ăn.}

\vocabentry{Knead}
{/niːd/ (v)}
{Work (moistened flour or clay) into dough or paste with the hands.}
{You need to knead the dough for at least ten minutes to make bread.}
{Đây là một động từ dùng để chỉ việc nhào bột.}

\vocabentry{Leavening}
{/ˈlevənɪŋ/ (n)}
{A substance used in dough or batter to make it rise.}
{Yeast is the most common leavening agent used in traditional baking.}
{Đây là một danh từ dùng để chỉ chất gây nở, ví dụ như men hoặc bột nở.}

\vocabentry{Marinate}
{/ˈmærɪneɪt/ (v)}
{Soak (meat, fish, or other food) in a marinade.}
{It is best to marinate the chicken overnight for maximum flavor.}
{Đây là một động từ dùng để chỉ việc tẩm ướp thức ăn.}

\vocabentry{Nutritious}
{/nuˈtrɪʃəs/ (adj)}
{Efficient as food; nourishing.}
{Eggs are a highly nutritious food that is high in protein and vitamins.}
{Đây là một tính từ dùng để chỉ thực phẩm giàu dinh dưỡng và bổ dưỡng.}

\vocabentry{Pasteurize}
{/ˈpæstʃəraɪz/ (v)}
{Subject (milk, wine, or other food) to a process of partial sterilization.}
{Most milk sold in supermarkets has been pasteurized to kill harmful bacteria.}
{Đây là một động từ dùng để chỉ việc tiệt trùng thực phẩm.}

\vocabentry{Quench}
{/kwentʃ/ (v)}
{Satisfy (one's thirst) by drinking.}
{A cold glass of lemonade is perfect for quenching your thirst on a hot day.}
{Đây là một động từ dùng để chỉ việc giải khát hoặc làm dịu cơn khát.}

\vocabentry{Rancid}
{/ˈrænsɪd/ (adj)}
{(Of foods containing fat or oil) smelling or tasting unpleasant as a result of being old and stale.}
{The butter had gone rancid because it wasn't kept in the refrigerator.}
{Đây là một tính từ dùng để chỉ thực phẩm bị ôi thiu, khét mùi dầu mỡ.}

\vocabentry{Savor}
{/ˈseɪvər/ (v)}
{Taste (good food or drink) and enjoy it completely.}
{He closed his eyes to savor the delicious flavor of the chocolate cake.}
{Đây là một động từ dùng để chỉ việc thưởng thức hoặc nhấm nháp hương vị thức ăn.}

\vocabentry{Texture}
{/ˈtekstʃər/ (n)}
{The feel, appearance, or consistency of a surface or a substance.}
{The creamy texture of the soup made it very satisfying and comforting.}
{Đây là một danh từ dùng để chỉ kết cấu của thức ăn như độ giòn, mềm, dai.}

\vocabentry{Unprocessed}
{/ˌʌnˈproʊsest/ (adj)}
{Not having been through any industrial treatment or process.}
{A diet based on unprocessed foods is generally much healthier for the body.}
{Đây là một tính từ dùng để chỉ thực phẩm thô, chưa qua chế biến.}

\vocabentry{Variety}
{/vəˈraɪəti/ (n)}
{The quality or state of being different or diverse.}
{Eating a wide variety of fruits and vegetables ensures you get all essential nutrients.}
{Đây là một danh từ dùng để chỉ sự đa dạng trong khẩu phần ăn.}

\vocabentry{Wholesome}
{/ˈhoʊlsəm/ (adj)}
{Conducive to or suggestive of good health and physical well-being.}
{The farmhouse provided a wholesome meal of fresh bread, cheese, and milk.}
{Đây là một tính từ dùng để chỉ thực phẩm lành mạnh và tốt cho sức khỏe.}

\vocabentry{Yeast}
{/jiːst/ (n)}
{A microscopic fungus consisting of single oval cells that reproduce by budding.}
{Yeast is essential for making the dough rise when baking bread.}
{Đây là một danh từ dùng để chỉ men làm bánh.}

\vocabentry{Zest}
{/zest/ (n)}
{The outer colored part of the peel of a citrus fruit, used as flavoring.}
{Add a little lemon zest to the cake batter for a fresh and citrusy flavor.}
{Đây là một danh từ dùng để chỉ vỏ quả thuộc họ cam chanh dùng làm hương liệu.}

\vocabentry{Appetite}
{/ˈæpɪtaɪt/ (n)}
{A natural desire to satisfy a bodily need, especially for food.}
{Fresh air and exercise can help to increase your appetite for dinner.}
{Đây là một danh từ dùng để chỉ sự thèm ăn.}

\vocabentry{Balance}
{/ˈbæləns/ (v)}
{A condition in which different elements are equal or in the correct proportions.}
{You should balance your intake of proteins, fats, and carbohydrates.}
{Sự cân bằng sinh thái.}

\vocabentry{Confectionery}
{/kənˈfekʃəneri/ (n)}
{Sweets and chocolates considered collectively.}
{The shop specializes in traditional handmade confectionery and chocolates.}
{Đây là một danh từ dùng để chỉ các loại bánh kẹo đồ ngọt nói chung.}

\vocabentry{Dietary}
{/ˈdaɪəteri/ (adj)}
{Relating to the transition of food or a person's diet.}
{Fiber is an important dietary component that aids in proper digestion.}
{Đây là một tính từ dùng để chỉ những gì thuộc về chế độ ăn uống.}

\vocabentry{Emulsifier}
{/ɪˈmʌlsɪfaɪər/ (n)}
{A substance that stabilizes an emulsion.}
{Lecithin is a common emulsifier used in the production of chocolate.}
{Đây là một danh từ dùng để chỉ chất nhũ hóa trong thực phẩm.}

\vocabentry{Fortify}
{/ˈfɔːrtɪfaɪ/ (v)}
{Strengthen (food) with vitamins or minerals.}
{Many breakfast cereals are fortified with iron and other essential vitamins.}
{Đây là một động từ dùng để chỉ việc bổ sung thêm vi chất vào thực phẩm.}

\vocabentry{Glucose}
{/ˈɡluːkoʊs/ (n)}
{A simple sugar which is an important energy source in living organisms.}
{The body converts carbohydrates into glucose to provide energy for the cells.}
{Đây là một danh từ dùng để chỉ đường glucose.}

\vocabentry{Hydroponic}
{/ˌhaɪdrəˈpɑːnɪk/ (adj)}
{Relating to the process of growing plants in sand, gravel, or liquid, with added nutrients but without soil.}
{Hydroponic farming allows for the production of fresh vegetables in urban areas.}
{Đây là một tính từ dùng để chỉ phương pháp trồng cây thủy canh.}

\vocabentry{Intolerance}
{/ɪnˈtɑːlərəns/ (n)}
{An inability to eat a food or take a drug without adverse effects.}
{Lactose intolerance makes it difficult for some people to digest dairy products.}
{Đây là một danh từ dùng để chỉ sự không dung nạp thức ăn.}

\vocabentry{Juice}
{/dʒuːs/ (n/v)}
{The liquid parts of fruits or vegetables; to extract the liquid parts.}
{Fresh orange juice is a popular and healthy way to start the day.}
{Đây là một danh từ/động từ dùng để chỉ nước ép hoặc việc ép nước trái cây.}

\vocabentry{Kernel}
{/ˈkɜːrnl/ (n)}
{A softer, usually edible part of a nut, seed, or fruit stone contained within its shell.}
{The walnuts are prized for their large and delicious kernels.}
{Đây là một danh từ dùng để chỉ nhân hoặc hạt của các loại quả hạch.}

\vocabentry{Lactose}
{/ˈlæktoʊs/ (n)}
{A sugar present in milk.}
{Some people lack the enzyme needed to break down lactose in dairy food.}
{Đây là một danh từ dùng để chỉ đường sữa.}

\vocabentry{Munch}
{/mʌntʃ/ (v)}
{Eat (something) with a continuous and often audible action of the jaws.}
{She spent the whole afternoon munching on apples and crackers.}
{Đây là một động từ dùng để chỉ việc nhai tóp tép hoặc ăn vặt.}

\vocabentry{Nourish}
{/ˈnɜːrɪʃ/ (v)}
{Provide with the food or other substances necessary for growth, health, and good condition.}
{Healthy food is essential to nourish the body and keep it strong.}
{Đây là một động từ dùng để chỉ việc nuôi dưỡng cơ thể.}

\vocabentry{Overindulge}
{/ˌoʊvərɪnˈdʌldʒ/ (v)}
{Have too much of something enjoyable, especially food or drink.}
{It is easy to overindulge in sweets and desserts during the holiday season.}
{Đây là một động từ dùng để chỉ việc ăn uống quá độ hoặc nuông chiều bản thân quá mức.}

\vocabentry{Pungent}
{/ˈpʌndʒənt/ (adj)}
{Having a sharply strong taste or smell.}
{Garlic and onions have a very pungent smell that can be quite overpowering.}
{Đây là một tính từ dùng để chỉ mùi vị hăng hoặc cay nồng của thực phẩm.}

\vocabentry{Quantity}
{/ˈkwɑːntəti/ (n)}
{The amount or number of a material or immaterial thing.}
{The restaurant serves high-quality food in very generous quantities.}
{Đây là một danh từ dùng để chỉ số lượng thức ăn.}

\vocabentry{Refined}
{/rɪˈfaɪnd/ (adj)}
{(Of a substance) having been processed to remove impurities or unwanted elements.}
{Reducing your intake of refined sugars can significantly improve your health.}
{Đây là một tính từ dùng để chỉ thực phẩm đã qua tinh chế như đường trắng, bột mì trắng.}

\vocabentry{Saturated}
{/ˈsætʃəreɪtɪd/ (adj)}
{(Of a fat) containing a high proportion of fatty acid molecules without double bonds.}
{Diets high in saturated fats are linked to an increased risk of heart disease.}
{Đây là một tính từ dùng để chỉ chất béo bão hòa, thường không tốt cho tim mạch.}

\vocabentry{Tasteless}
{/ˈteɪstləs/ (adj)}
{Lacking flavor.}
{I found the vegetables to be rather tasteless because they were overcooked.}
{Đây là một tính từ dùng để chỉ thực phẩm vô vị hoặc không có mùi vị.}

\vocabentry{Unsavoury}
{/ʌnˈseɪvəri/ (adj)}
{Unpleasant to taste, smell, or look at.}
{The meat had an unsavoury smell, so I decided not to eat it.}
{Đây là một tính từ dùng để chỉ thực phẩm không ngon hoặc có mùi vị khó chịu.}

\vocabentry{Vessel}
{/ˈvesl/ (n)}
{A container in which liquid may be kept.}
{Use a clean glass vessel to store the homemade salad dressing.}
{Trong ẩm thực, đây là danh từ chỉ dụng cụ đựng thực phẩm hoặc bình chứa.}

\vocabentry{Wok}
{/wɑːk/ (n)}
{A bowl-shaped frying pan used typically in Chinese cooking.}
{Stir-frying in a hot wok is a quick and healthy way to cook vegetables.}
{Đây là một danh từ dùng để chỉ cái chảo sâu lòng kiểu Á.}

\vocabentry{Yield}
{/jiːld/ (n/v)}
{An amount produced of an agricultural or industrial product; to produce.}
{This recipe will yield enough pasta to serve four hungry people.}
{Đây là một danh từ/động từ chỉ sản lượng thực phẩm hoặc việc mang lại kết quả.}

\vocabentry{Zesty}
{/ˈzesti/ (adj)}
{Having a strong, pleasant, and somewhat spicy flavor. (Đây là một tính từ dùng để chỉ hương vị nồng nàn, tươi mát (thường là vị cam chanh).).}
{The salad dressing has a zesty flavor thanks to the addition of lime juice.}
{}

\vocabentry{Aerobic}
{/erˈoʊbɪk/ (adj)}
{Relating to, involving, or requiring free oxygen.}
{Aerobic metabolism uses oxygen to convert nutrients into energy for muscles.}
{Trong dinh dưỡng, đây là tính từ chỉ quá trình chuyển hóa hiếu khí.}

\vocabentry{Beverage}
{/ˈbevərɪdʒ/ (n)}
{A drink, especially one other than water.}
{We offer a wide selection of hot and cold beverages in our cafe.}
{Đây là một danh từ trang trọng dùng để chỉ các loại đồ uống.}

\vocabentry{Condiment}
{/ˈkɑːndɪmənt/ (n)}
{A substance such as salt or mustard that is used to add flavor to food.}
{Ketchup and mayonnaise are two of the most popular condiments in the world.}
{Đây là một danh từ dùng để chỉ đồ gia vị ăn kèm.}

\vocabentry{Dehydrate}
{/ˌdiːˈhaɪdreɪt/ (v)}
{Cause (a person or their body) to lose a large amount of water.}
{You can dehydrate fruit to make healthy and delicious snacks for traveling.}
{Đây là một động từ dùng để chỉ việc làm mất nước hoặc làm khô thực phẩm.}

\vocabentry{Enrich}
{/ɪnˈrɪtʃ/ (v)}
{Improve or enhance the quality or value of.}
{White bread is often enriched with vitamins that were lost during processing.}
{Trong thực phẩm, đây là việc làm giàu thêm dưỡng chất.}

\vocabentry{Flavouring}
{/ˈfleɪvərɪŋ/ (n)}
{A substance used to give a particular flavor to food or drink.}
{Vanilla is a common flavouring used in many cakes and desserts.}
{Đây là một danh từ dùng để chỉ hương liệu thực phẩm.}

\vocabentry{Gluttony}
{/ˈɡlʌtəni/ (n)}
{Habitual greed or excess in eating.}
{Gluttony is considered one of the seven deadly sins in some cultures.}
{Đây là một danh từ dùng để chỉ thói ham ăn quá độ.}

\vocabentry{Hearth}
{/hɜːrθ/ (n)}
{The floor of a fireplace.}
{The bread was baked in a traditional oven over an open hearth.}
{Trong ẩm thực truyền thống, đây là danh từ chỉ bếp lửa.}

\vocabentry{Inedible}
{/ɪnˈedəbl/ (adj)}
{Not fit or safe to be eaten.}
{The fruit was completely inedible because it was covered in mold.}
{Đây là một tính từ dùng để chỉ thực phẩm không thể ăn được.}

\vocabentry{Jelly}
{/ˈdʒeli/ (n)}
{A sweet clear food made from fruit juice and sugar boiled together.}
{She spread a thin layer of strawberry jelly on her toast for breakfast.}
{Đây là một danh từ dùng để chỉ thạch hoặc mứt dẻo.}

\vocabentry{Kilojoule}
{/ˈkɪlədʒuːl/ (n)}
{A unit of energy equal to 1,000 joules. (Đây là một danh từ dùng để chỉ đơn vị năng lượng thực phẩm (kJ).).}
{The nutritional label shows the number of kilojoules in each serving of the product.}
{}

\vocabentry{Legume}
{/ˈleɡjuːm/ (n)}
{A leguminous plant, especially one grown as a crop.}
{Lentils and chickpeas are nutritious legumes that are high in protein and fiber.}
{Đây là một danh từ dùng để chỉ các loại cây họ đậu.}

\vocabentry{Nutrient-dense}
{/ˈnuːtriənt dens/ (adj)}
{Rich in nutrients relative to their calorie content.}
{Leafy greens are nutrient-dense foods that provide many health benefits.}
{Đây là một tính từ dùng để chỉ thực phẩm có mật độ dinh dưỡng cao.}

\vocabentry{Oven}
{/ˈʌvn/ (n)}
{An enclosed compartment used for baking, heating, or drying food.}
{Preheat the oven to 180 degrees before putting the cake inside.}
{Đây là một danh từ dùng để chỉ lò nướng.}

\vocabentry{Protein-rich}
{/ˈproʊtiːn rɪtʃ/ (adj)}
{Containing a large amount of protein.}
{A protein-rich diet is important for athletes who want to build muscle mass.}
{Đây là một tính từ dùng để chỉ thực phẩm giàu chất đạm.}

\vocabentry{Quality}
{/ˈkwɑːləti/ (n)}
{The standard of something as measured against other things; a distinctive attribute.}
{Food quality is a major concern for consumers who shop at organic markets.}
{Đây là một danh từ dùng để chỉ chất lượng/phẩm chất.}

\vocabentry{Ripen}
{/ˈraɪpən/ (v)}
{(Of fruit or grain) become or make ripe.}
{You can leave the bananas on the counter to ripen for a few days.}
{Đây là một động từ dùng để chỉ việc làm chín hoặc quả chín.}

\vocabentry{Succulent}
{/ˈsʌkjələnt/ (adj)}
{(Of food) tender, juicy, and tasty.}
{The grilled steak was incredibly succulent and full of flavor.}
{Đây là một tính từ dùng để chỉ thực phẩm mọng nước và ngon lành.}

\vocabentry{Trace element}
{/ˈtreɪs ˌelɪmənt/ (n)}
{A chemical element present only in minute amounts.}
{Iron and zinc are essential trace elements that the body needs in small amounts.}
{Đây là một danh từ dùng để chỉ nguyên tố vi lượng trong dinh dưỡng.}

\vocabentry{Unsaturated}
{/ˌʌnˈsætʃəreɪtɪd/ (adj)}
{(Of a fat) containing double bonds between some of the carbon atoms.}
{Olive oil is a great source of healthy unsaturated fats.}
{Đây là một tính từ dùng để chỉ chất béo không bão hòa, thường tốt cho sức khỏe.}

\vocabentry{Vitamin-rich}
{/ˈvaɪtəmɪn rɪtʃ/ (adj)}
{Containing a large amount of vitamins.}
{Fresh berries are vitamin-rich snacks that can improve your overall health.}
{Đây là một tính từ dùng để chỉ thực phẩm giàu vitamin.}

\vocabentry{Whey}
{/weɪ/ (n)}
{The watery part of milk that remains after the formation of curds.}
{Whey protein is a popular supplement for people who exercise regularly.}
{Đây là một danh từ dùng để chỉ váng sữa hoặc đạm whey.}

\vocabentry{Yolk}
{/joʊk/ (n)}
{The yellow internal part of a bird's egg.}
{The recipe calls for three egg yolks to make the custard thick and creamy.}
{Đây là một danh từ dùng để chỉ lòng đỏ trứng.}

\vocabentry{Zinc}
{/zɪŋk/ (n)}
{A chemical element that is an essential trace element for humans.}
{Zinc plays a vital role in supporting a healthy immune system.}
{Đây là một danh từ dùng để chỉ kẽm, một vi chất quan trọng.}

\vocabentry{Absorption}
{/əbˈsɔːrpʃn/ (n)}
{The process or action by which one thing absorbs or is absorbed by another.}
{Vitamin D is necessary for the proper absorption of calcium in the body.}
{Trong dinh dưỡng, đây là sự hấp thụ dưỡng chất.}

\vocabentry{Conserve}
{/kənˈsɜːrv/ (v)}
{Protect (something, especially an environmentally or culturally important place or thing) from harm or destruction.}
{Steaming vegetables is a great way to conserve their natural nutrients and vitamins.}
{Trong nấu ăn là bảo toàn dưỡng chất.}

\vocabentry{Deficiency}
{/dɪˈfɪʃnsi/ (n)}
{A lack or shortage of something.}
{Vitamin A deficiency can lead to serious problems with vision and immunity.}
{Trong dinh dưỡng là sự thiếu hụt vi chất.}

\vocabentry{Essential}
{/ɪˈsenʃl/ (adj)}
{Absolutely necessary; extremely important.}
{Certain amino acids are essential because the body cannot produce them.}
{Đây là một tính từ dùng để chỉ sự thiết yếu.}

\vocabentry{Fats}
{/fæts/ (n)}
{A natural oily or greasy substance occurring in animal bodies, especially when deposited as a layer under the skin.}
{Not all fats are bad for you; some are essential for brain health and energy.}
{Đây là một danh từ dùng để chỉ chất béo.}

\vocabentry{Grains}
{/ɡreɪnz/ (n)}
{Wheat or any other cultivated cereal crop used as food.}
{Whole grains are much healthier than refined ones because they contain more fiber.}
{Đây là một danh từ dùng để chỉ ngũ cốc.}

\vocabentry{Herb}
{/hɜːrb/ (n)}
{Any plant with leaves, seeds, or flowers used for flavoring, food, medicine, or perfume.}
{I like to use fresh herbs like basil and rosemary to flavor my cooking.}
{Đây là một danh từ dùng để chỉ thảo mộc hoặc rau thơm.}

\vocabentry{Inhale}
{/ɪnˈheɪl/ (v)}
{Breathe in (air, gas, smoke, etc.).}
{He stopped to inhale the wonderful aroma of fresh bread from the bakery.}
{Trong ẩm thực là hít hà hương thơm.}

\vocabentry{Jar}
{/dʒɑːr/ (n)}
{A wide-mouthed cylindrical container made of glass or pottery.}
{I keep my homemade jam in small glass jars in the refrigerator.}
{Đây là một danh từ dùng để chỉ hũ hoặc lọ đựng thực phẩm.}

\vocabentry{Kilocalorie}
{/ˈkɪləkæləri/ (n)}
{A unit of energy equal to 1,000 calories.}
{Most adults need between 2,000 and 2,500 kilocalories a day to maintain their weight.}
{Đơn vị đo năng lượng thực phẩm phổ biến nhất.}

\vocabentry{Liquid}
{/ˈlɪkwɪd/ (n/adj)}
{A substance that flows freely but is of constant volume.}
{Soups and smoothies are good ways to consume nutrients in liquid form.}
{Trong dinh dưỡng chỉ thực phẩm dạng lỏng.}

\vocabentry{Maintain}
{/meɪnˈteɪn/ (v)}
{Cause or enable (a condition or state of affairs) to continue.}
{A healthy diet is essential to maintain a stable weight and good health.}
{Trong sức khỏe là duy trì thể trạng.}

\vocabentry{Natural}
{/ˈnætʃrəl/ (adj)}
{Existing in or caused by nature; not made or caused by humankind.}
{Natural sweeteners like honey are often better for you than artificial ones.}
{Đây là một tính từ dùng để chỉ thực phẩm tự nhiên.}

\vocabentry{Optimal}
{/ˈɑːptɪməl/ (adj)}
{Best or most favorable.}
{For optimal health, you should drink plenty of water and get enough sleep.}
{Trong dinh dưỡng chỉ mức độ tối ưu.}

\vocabentry{Preserve}
{/prɪˈzɜːrv/ (v)}
{Maintain (something) in its original or existing state.}
{Drying and salting are two ancient methods used to preserve meat for long periods.}
{Trong thực phẩm là bảo quản thức ăn.}

\vocabentry{Quick}
{/kwɪk/ (adj)}
{Moving fast or doing something in a short time.}
{A stir-fry is a quick and healthy meal that you can make in under twenty minutes.}
{Trong ẩm thực là các món ăn nhanh hoặc sơ chế nhanh.}

\vocabentry{Reduce}
{/rɪˈduːs/ (v)}
{Make smaller or less in amount, degree, or size.}
{You should try to reduce your intake of salt and saturated fats.}
{Trong dinh dưỡng là giảm lượng tiêu thụ.}

\vocabentry{Substance}
{/ˈsʌbstəns/ (n)}
{A particular kind of matter with uniform properties.}
{Milk is a complex substance containing proteins, fats, sugars, and vitamins.}
{Chất hoặc thành phần có trong thực phẩm.}

\vocabentry{Trace}
{/treɪs/ (n)}
{A very small amount of something.}
{The product may contain traces of nuts, so people with allergies should be careful.}
{Lượng vết hoặc lượng cực nhỏ.}

\vocabentry{Ultimate}
{/ˈʌltɪmət/ (adj)}
{Being or happening at the end of a process; final.}
{The ultimate goal of a healthy diet is to prevent disease and promote long life.}
{Mục tiêu cuối cùng của dinh dưỡng.}

\vocabentry{Valuable}
{/ˈvæljuəbl/ (adj)}
{Worth a great deal of money.}
{Oily fish are a valuable source of omega-3 fatty acids for the brain.}
{Trong dinh dưỡng là thực phẩm có giá trị cao về mặt sức khỏe.}

\vocabentry{Wholesale}
{/ˈhoʊlseɪl/ (adj/adv)}
{Done on a large scale; extensive.}
{Buying food in wholesale quantities can help you save a significant amount of money.}
{Trong thực phẩm chỉ việc mua bán sỉ.}

\vocabentry{Xylitol}
{/ˈzaɪlɪtɔːl/ (n)}
{A sweet-tasting crystalline alcohol derived from xylose.}
{Xylitol is often used as a sugar substitute in sugar-free chewing gum.}
{Đây là một danh từ dùng để chỉ một loại đường nhân tạo ít calo.}

\vocabentry{Yielding}
{/ˈjiːldɪŋ/ (adj)}
{(Of a substance) giving way under pressure.}
{The ripe peach was soft and yielding to the touch.}
{Trong ẩm thực chỉ thực phẩm mềm hoặc dễ nhai.}

\vocabentry{Zone}
{/zoʊn/ (n)}
{An area or stretch of land having a particular characteristic.}
{People living in Blue Zones have diets that are very high in plant-based foods.}
{Trong dinh dưỡng thường dùng cho các "vùng xanh" - Blue Zones nơi thọ nhất thế giới.}

\vocabentry{Adulteration}
{/əˌdʌltəˈreɪʃn/ (n)}
{The action of making something poorer in quality by the addition of another substance.}
{Food adulteration is a serious crime that can pose a major risk to public health.}
{Sự làm giả hoặc pha tạp thực phẩm.}

\vocabentry{Balanced}
{/ˈbælənst/ (adj)}
{Taking everything into account; fairly judged or presented.}
{A balanced meal should include protein, carbohydrates, and healthy fats.}
{Tính từ chỉ sự cân đối của chế độ ăn.}

\vocabentry{Consumption}
{/kənˈsʌmpʃn/ (n)}
{The using up of a resource.}
{Global meat consumption has increased significantly over the last few decades.}
{Sự tiêu thụ thực phẩm.}

\vocabentry{Deficient}
{/dɪˈfɪʃnt/ (adj)}
{Not having enough of a specified quality or ingredient.}
{If your diet is deficient in Vitamin C, you might develop scurvy.}
{Tính từ chỉ tình trạng bị thiếu hụt vi chất.}

\vocabentry{Excessive}
{/ɪkˈsesɪv/ (adj)}
{More than is necessary, normal, or desirable; immoderate.}
{Excessive consumption of sugar can lead to weight gain and tooth decay.}
{Tính từ chỉ sự quá mức, dư thừa.}

\vocabentry{Function}
{/ˈfʌŋkʃn/ (n)}
{An activity or purpose natural to or intended for a person or thing.}
{One of the main functions of protein is to build and repair body tissues.}
{Chức năng của các chất trong cơ thể.}

\vocabentry{Greed}
{/ɡriːd/ (n)}
{Intense and selfish desire for something, especially wealth, power, or food.}
{Obesity is often caused by a combination of genetics and individual greed for food.}
{Sự tham ăn hoặc tham lam thực phẩm.}

\vocabentry{Heaviness}
{/ˈhevinəs/ (n)}
{The quality of being heavy.}
{I felt a sense of heaviness in my stomach after eating that large pizza.}
{Cảm giác nặng bụng sau khi ăn quá nhiều.}

\vocabentry{Intake}
{/ˈɪnteɪk/ (n)}
{An amount of food, air, or another substance taken into the body.}
{You should monitor your daily intake of calories if you are trying to lose weight.}
{Lượng thực phẩm đưa vào cơ thể.}

\vocabentry{Joint}
{/dʒɔɪnt/ (n)}
{A structure in the human or animal body at which two parts of the skeleton are fitted together.}
{Nutrients like glucosamine are important for maintaining healthy joints.}
{Xương khớp - chịu ảnh hưởng của dinh dưỡng như Canxi.}

\vocabentry{Kilogram}
{/ˈkɪləɡræm/ (n)}
{The SI unit of mass.}
{He lost five kilograms after following a strict diet and exercise plan for a month.}
{Đơn vị đo khối lượng thực phẩm hoặc cơ thể.}

\vocabentry{Life}
{/laɪf/ (n)}
{The condition that distinguishes animals and plants from inorganic matter.}
{Access to nutritious food is a fundamental human right and essential for life.}
{Dinh dưỡng là cội nguồn của sự sống.}

\vocabentry{Measurement}
{/ˈmeʒərmənt/ (n)}
{The action of measuring something.}
{Accurate measurement of ingredients is essential for successful baking.}
{Sự đo lường calo hoặc thành phần dinh dưỡng.}

\vocabentry{Necessity}
{/nəˈsesəti/ (n)}
{The fact of being required or indispensable.}
{Food and water are the most basic necessities for human survival.}
{Thực phẩm là nhu yếu phẩm thiết yếu.}

\vocabentry{Order}
{/ˈɔːrdər/ (n/v)}
{The arrangement or disposition of people or things in relation to each other.}
{I would like to order a large salad and a bottle of mineral water, please.}
{Trong ẩm thực chỉ việc đặt món ăn.}

\vocabentry{Physical}
{/ˈfɪzɪkl/ (adj)}
{Relating to the body as opposed to the mind.}
{Proper nutrition is essential for maintaining physical strength and energy.}
{Thể trạng vật lý chịu tác động bởi dinh dưỡng.}

\vocabentry{Range}
{/reɪndʒ/ (n)}
{A set of different things of the same general type.}
{The supermarket offers a wide range of organic and health food products.}
{Dải thực phẩm hoặc phạm vi dưỡng chất.}

\vocabentry{Standard}
{/ˈstændərd/ (n)}
{A level of quality or attainment.}
{The government has strict standards for the labeling of nutritional information on food.}
{Tiêu chuẩn an toàn thực phẩm.}

\vocabentry{Topic}
{/ˈtɑːpɪk/ (n)}
{A matter dealt with in a text, discourse, or conversation.}
{Health and nutrition is a very popular topic for the IELTS Speaking test.}
{Đề tài dinh dưỡng phổ biến trong IELTS.}

\vocabentry{Underweight}
{/ˌʌndərˈweɪt/ (adj)}
{Below a weight considered normal or desirable.}
{Being severely underweight can lead to problems with immunity and bone health.}
{Tính từ chỉ tình trạng thiếu cân.}

\vocabentry{Vigor}
{/ˈvɪɡər/ (n)}
{Physical strength and good health.}
{A healthy diet can help you regain your vigor and energy after an illness.}
{Sức lực và sinh khí do ăn uống tốt mang lại.}

\vocabentry{Waste}
{/weɪst/ (n/v)}
{Material which is eliminated or discarded as no longer useful or required.}
{We should all try to reduce the amount of food waste we produce each week.}
{Rác thải thực phẩm hoặc sự lãng phí thức ăn.}

\vocabentry{Zero}
{/ˈzɪroʊ/ (n/adj)}
{No quantity or number.}
{Diet sodas are often marketed as having zero calories and zero sugar.}
{Trong thực phẩm là "zero calories" - không calo.}

\vocabentry{Adequate}
{/ˈædɪkwət/ (adj)}
{Satisfactory or acceptable in quality or quantity.}
{Many people in the world still do not have an adequate supply of nutritious food.}
{Tính từ chỉ sự đầy đủ dưỡng chất.}

\vocabentry{Bulk}
{/bʌlk/ (n)}
{The mass or size of something large.}
{Buying food in bulk can be a great way to save money and reduce packaging waste.}
{Thực phẩm thô hoặc mua số lượng lớn.}

\vocabentry{Complex}
{/ˈkɑːmpleks/ (adj)}
{Consisting of many different and connected parts.}
{Complex carbohydrates like brown rice provide a steady release of energy.}
{Trong dinh dưỡng là carbohydrate phức hợp - tốt cho sức khỏe.}

\vocabentry{Diverse}
{/daɪˈvɜːrs/ (adj)}
{Showing a great deal of variety; very different.}
{The city has a very diverse food scene with restaurants from all over the world.}
{Tính từ chỉ sự đa dạng trong ẩm thực.}

\vocabentry{Fortification}
{/ˌfɔːrtɪfɪˈkeɪʃn/ (n)}
{The action of fortifying someone or something.}
{The fortification of flour with folic acid has helped to reduce birth defects.}
{Sự bổ sung vi chất vào thực phẩm.}

\vocabentry{Grow}
{/ɡroʊ/ (v)}
{(Of a living thing) undergo natural development by increasing in size.}
{Children need a lot of protein and calcium to grow tall and strong.}
{Sự tăng trưởng của cơ thể hoặc cây trồng thực phẩm.}

\vocabentry{High}
{/haɪ/ (adj)}
{Of great vertical extent.}
{You should avoid snacks that are high in sugar and saturated fats.}
{Trong dinh dưỡng là thực phẩm "high in..." - giàu cái gì đó.}

\vocabentry{Iron}
{/ˈaɪərn/ (n)}
{A chemical element that is an essential mineral for the body.}
{Spinach and red meat are excellent sources of dietary iron.}
{Sắt - chất quan trọng cho máu.}

\vocabentry{Juiciness}
{/ˈdʒuːsinəs/ (n)}
{The quality of being juicy.}
{The juiciness of the peach made it very refreshing on a hot day.}
{Độ mọng nước của hoa quả hoặc thịt.}

\vocabentry{Knowledgeable}
{/ˈnɑːlɪdʒəbl/ (adj)}
{Intelligent and well informed.}
{Being knowledgeable about nutrition can help you make healthier food choices.}
{Sự am hiểu về kiến thức dinh dưỡng.}

\vocabentry{Long-term}
{/ˌlɔːŋ ˈtɜːrm/ (adj)}
{Occurring over or relating to a long period of time.}
{Poor eating habits in childhood can have serious long-term consequences for health.}
{Hệ quả dài hạn của thói quen ăn uống.}

\vocabentry{Modest}
{/ˈmɑːdɪst/ (adj)}
{(Of an amount, rate, or level) relatively moderate, limited, or small.}
{Eating a modest amount of dark chocolate can actually be good for your heart.}
{Trong dinh dưỡng là khẩu phần khiêm tốn/vừa phải.}

\vocabentry{Nutritional}
{/nuˈtrɪʃənl/ (adj)}
{Relating to nutrition.}
{You should always check the nutritional label before buying processed food.}
{Tính từ chỉ những gì thuộc về giá trị dinh dưỡng.}

\vocabentry{Overweight}
{/ˌoʊvərˈweɪt/ (adj)}
{Above a weight considered normal or desirable.}
{Many children in developed nations are now considered to be overweight.}
{Tính từ chỉ tình trạng thừa cân.}

\vocabentry{Pure}
{/pjʊr/ (adj)}
{Not mixed or adulterated with any other substance or material.}
{This honey is 100\% pure and has no added sugar or preservatives.}
{Tính từ chỉ thực phẩm nguyên chất, tinh khiết.}

\vocabentry{Reliable}
{/rɪˈlaɪəbl/ (adj)}
{Consistently good in quality or performance; able to be trusted.}
{It is important to find a reliable source of information about healthy eating.}
{Trong dinh dưỡng là nguồn tin hoặc thực phẩm đáng tin.}

\vocabentry{Sustainable}
{/səˈsteɪnəbl/ (adj)}
{Able to be maintained at a certain rate or level.}
{Sustainable agriculture aims to produce food without damaging the environment.}
{Đây là một tính từ dùng để chỉ sự bền vững về năng lượng và môi trường.}

\vocabentry{Traditional}
{/trəˈdɪʃənl/ (adj)}
{Existing in or as part of a tradition; long-established.}
{The festival features a wide variety of traditional local dishes.}
{Ẩm thực truyền thống.}

\vocabentry{Underestimate}
{/ˌʌndərˈestɪmeɪt/ (v)}
{Estimate (something) to be smaller or less important than it actually is.}
{Many people underestimate the number of calories in a large soda.}
{Đánh giá thấp calo hoặc tác hại của đồ ăn nhanh.}

\vocabentry{Vital}
{/ˈvaɪtl/ (adj)}
{Absolutely necessary or important; essential.}
{Water is vital for every single process that happens in the human body.}
{Cực kỳ quan trọng đối với sức khỏe.}

\vocabentry{Weakness}
{/ˈwiːknəs/ (n)}
{The state or condition of lacking strength.}
{Iron deficiency can cause a persistent sense of fatigue and physical weakness.}
{Sự suy nhược do thiếu dinh dưỡng.}

\vocabentry{Yellow}
{/ˈjeloʊ/ (adj)}
{Of the color between green and orange in the spectrum.}
{Yellow fruits and vegetables like mangoes and carrots are high in Vitamin A.}
{Màu sắc chỉ thực phẩm giàu beta-carotene.}

\vocabentry{Zeal}
{/ziːl/ (n)}
{Great energy or enthusiasm in pursuit of a cause or an objective.}
{She approached her new diet with great zeal and dedication.}
{Lòng nhiệt huyết với việc ăn uống lành mạnh.}

\vocabentry{Amount}
{/əˈmaʊnt/ (n)}
{A quantity of something, especially the total of a thing or things in number, size, value, or extent.}
{The amount of sugar in some fruit juices can be surprisingly high.}
{Số lượng vi chất hoặc calo.}

\vocabentry{Benefit}
{/ˈbenɪfɪt/ (n/v)}
{An advantage or profit gained from something.}
{Regular consumption of oily fish has many benefits for brain and heart health.}
{Lợi ích của thực phẩm tốt.}

\vocabentry{Cause}
{/kɔːz/ (n/v)}
{A person or thing that gives rise to an action, phenomenon, or condition.}
{Poor diet is a major cause of many preventable diseases in the world today.}
{Nguyên nhân của các bệnh liên quan đến ăn uống.}

\vocabentry{Description}
{/dɪˈskrɪpʃn/ (n)}
{A spoken or written representation or account of a person, object, or event.}
{The menu provides a detailed description of each dish and its ingredients.}
{Mô tả hương vị hoặc thành phần thức ăn.}

\vocabentry{Effect}
{/ɪˈfekt/ (n/v)}
{A change which is a result or consequence of an action or other cause.}
{The long-term effects of consuming high levels of caffeine are still being studied.}
{Tác động của thực phẩm lên cơ thể.}

\vocabentry{Frequency}
{/ˈfriːkwənsi/ (n)}
{The rate at which something occurs or is repeated over a particular period of time.}
{Increasing the frequency of your meals can help keep your metabolism steady.}
{Tần suất bữa ăn.}

\vocabentry{Growth}
{/ɡroʊθ/ (n)}
{The process of increasing in physical size.}
{Proper nutrition during pregnancy is essential for the healthy growth of the baby.}
{Sự trưởng thành cần dinh dưỡng.}

\vocabentry{Habit}
{/ˈhæbɪt/ (n)}
{A settled or regular tendency or practice, especially one that is hard to give up.}
{Developing healthy eating habits early in life can prevent many future problems.}
{Thói quen ăn uống.}

\vocabentry{Importance}
{/ɪmˈpɔːrtns/ (n)}
{The state or quality of being important.}
{Many nutritionists emphasize the importance of eating a healthy breakfast.}
{Tầm quan trọng của bữa sáng.}

\vocabentry{Justice}
{/ˈdʒʌstɪs/ (n)}
{Just behavior or treatment.}
{Food justice aims to ensure that everyone has access to healthy and affordable food.}
{Công bằng thực phẩm/an ninh lương thực.}

\vocabentry{Kind}
{/kaɪnd/ (n/adj)}
{A group of people or things having similar characteristics.}
{What kind of food do you usually prefer to eat when you go out?}
{Các loại thực phẩm.}

\vocabentry{Length}
{/leɪŋθ/ (n)}
{The measurement or extent of something from end to end.}
{The length of time you should cook meat depends on the type and size of the cut.}
{Độ dài của quá trình tiêu hóa.}

\vocabentry{Material}
{/məˈtɪriəl/ (n/adj)}
{The matter from which a thing is or can be made.}
{Raw material for food production must be kept in clean and cool conditions.}
{Nguyên liệu thô.}

\vocabentry{Number}
{/ˈnʌmbər/ (n/v)}
{An arithmetical value, expressed by a word, symbol, or figure.}
{The number of people suffering from food allergies has increased significantly.}
{Số lượng người mắc bệnh béo phì.}

\vocabentry{Opinion}
{/əˈpɪnjən/ (n)}
{A view or judgment formed about something.}
{In my opinion, the government should tax sugary drinks to improve public health.}
{Quan điểm về xu hướng ăn thuần chay.}

\vocabentry{Point}
{/pɔɪnt/ (n/v)}
{A particular spot, place, or moment in an enlarged whole or in a duration.}
{There is no point in taking supplements if your overall diet is very poor.}
{Thời điểm ăn uống.}

\vocabentry{Question}
{/ˈkwestʃən/ (n/v)}
{A sentence worded or expressed so as to elicit information.}
{The safety of genetically modified foods is still a major question for many people.}
{Vấn đề an toàn thực phẩm.}

\vocabentry{Reason}
{/ˈriːzn/ (n/v)}
{A cause, explanation, or justification for an action or event.}
{One reason many people choose organic food is to avoid chemical pesticides.}
{Lý do chọn thực phẩm hữu cơ.}

\vocabentry{Success}
{/səkˈses/ (n)}
{The accomplishment of an aim or purpose.}
{The success of your diet depends on your commitment and self-discipline.}
{Sự thành công của một chế độ ăn kiêng.}

\vocabentry{Thinking}
{/ˈθɪŋkɪŋ/ (n/adj)}
{The process of using one's mind to consider or reason about something.}
{Modern thinking about nutrition has shifted towards a more holistic approach.}
{Tư duy về thực phẩm sạch.}

\vocabentry{Unity}
{/ˈjuːnəti/ (n)}
{The state of being united or joined as a whole.}
{There is a lack of unity among experts regarding the benefits of certain diets.}
{Sự thống nhất trong tháp dinh dưỡng.}

\vocabentry{Volume}
{/ˈvɑːljuːm/ (n)}
{The amount of space that a substance or object occupies.}
{The volume of food you eat is just as important as its nutritional content.}
{Khối lượng thực phẩm tiêu thụ.}

\vocabentry{Year}
{/jɪər/ (n)}
{The time taken by the earth to make one revolution around the sun.}
{The national health report is published every year to track dietary trends.}
{Báo cáo dinh dưỡng hàng năm.}

\vocabentry{Zero-waste}
{/ˌzɪroʊ ˈweɪst/ (adj)}
{Producing little or no waste.}
{A zero-waste kitchen focuses on using every part of the ingredients to reduce trash.}
{Tính từ chỉ lối sống/nấu ăn không rác thải.}

\vocabentry{Acidity}
{/əˈsɪdəti/ (n)}
{The level of acid in substances such as water, soil, or food.}
{The high acidity of citrus fruits can sometimes cause stomach irritation.}
{Độ axit trong thực phẩm.}

\vocabentry{Baking}
{/ˈbeɪkɪŋ/ (n)}
{The activity of cooking food, especially cakes and bread, in an oven.}
{Baking is a precise science that requires accurate measurement of ingredients.}
{Nghệ thuật nướng bánh.}

\vocabentry{Calcium}
{/ˈkælsiəm/ (n)}
{A chemical element that is an essential mineral for healthy bones.}
{Dairy products and leafy greens are excellent sources of dietary calcium.}
{Canxi - khoáng chất thiết yếu.}

\vocabentry{Dough}
{/doʊ/ (n)}
{A thick, malleable mixture of flour and liquid, used for baking into bread or pastry.}
{You should let the dough rise in a warm place before putting it in the oven.}
{Bột nhào.}

\vocabentry{Enzyme}
{/ˈenzaɪm/ (n)}
{A substance produced by a living organism which acts as a catalyst.}
{Specific enzymes are needed to break down complex proteins into amino acids.}
{Men/Enzym tiêu hóa.}

\vocabentry{Fructose}
{/ˈfrʌktoʊs/ (n)}
{A sugar found especially in honey and fruit.}
{Fructose is a natural sugar that is often used as a sweetener in many products.}
{Đường hoa quả.}

\vocabentry{Gelatin}
{/ˈdʒelətn/ (n)}
{A virtually colorless and tasteless water-soluble protein prepared from collagen.}
{Gelatin is used to give structure and stability to desserts like panna cotta.}
{Chất làm đông thực phẩm.}

\vocabentry{Honeydew}
{/ˈhʌniduː/ (n)}
{A variety of melon with sweet green flesh.}
{Honeydew melon is a refreshing and hydrating snack during the summer months.}
{Dưa lưới xanh.}

\vocabentry{Infusion}
{/ɪnˈfjuːʒn/ (n)}
{A drink, remedy, or extract prepared by soaking the leaves of a plant or herb in liquid.}
{A warm herbal infusion can be very soothing after a large and heavy meal.}
{Nước hãm từ thảo mộc/trà.}

\vocabentry{Kebab}
{/kəˈbæb/ (n)}
{A dish of pieces of meat, fish, or vegetables roasted or grilled on a skewer.}
{Grilled chicken kebabs are a popular and healthy choice for a summer barbecue.}
{Món thịt nướng xiên.}

\vocabentry{Lipid}
{/ˈlɪpɪd/ (n)}
{Any of a class of organic compounds that are fatty acids.}
{Lipids are essential for cell membrane structure and the absorption of some vitamins.}
{Chất béo/Lipit.}

\vocabentry{Molasses}
{/məˈlæsɪz/ (n)}
{A thick, dark brown syrup obtained from raw sugar during the refining process.}
{Molasses is a rich source of iron and is often used in traditional baking recipes.}
{Mật rỉ đường.}

\vocabentry{Nitrate}
{/ˈnaɪtreɪt/ (n)}
{A salt or ester of nitric acid, often used as a preservative in cured meats.}
{High consumption of nitrates in processed meats has been linked to health risks.}
{Chất bảo quản muối nitrat.}

